% Detection, Forecasting, and Clinical Readiness (Merged from 07, 08, 09)
\subsection{Detection, Forecasting, and Clinical Readiness}
\label{sec:detection-forecasting-readiness}

Traditionally detection and forecasting detection serve largely different clinical purposes.

Detection systems alert caregivers when a seizure occurs enabling intervention during or immediately after the seizure. It is most valuable for generalized tonic-clonic or focal to bilateral tonic-clonic seizures where caregiver response can prevent injury or SUDEP.

In contrast, forecasting systems provide advance warning of increased seizure likelihood. This allows patients to modify activities, take rescue medication, or avoid risky situations. Forecasting could improve quality of life by reducing the uncertainty of when seizures might occur.

\subsubsection{Maturity Comparison}

Two detection studies achieve FPR performance comparable to Phase 3 studies in the ILAE systematic review. Several Commercial devices exist for detection including the Embrace watch with FDA 510(k) clearance and NightWatch with CE approval in Europe. Still no forecasting study meets a clear clinical benchmark. No device has regulatory approval for seizure forecasting.

While detection has a framework for classifying (ILAE 5 Phases by \textcite{beniczkyStandardsTestingClinical2018}), Forecasting lacks equivalent standardized guidelines. Forecasting AUC consistently falls between 0.74 and 0.80 across studies, suggesting a performance ceiling for non-invasive approaches, but the clinical acceptability of this performance level remains unclear.

Sample size disparities also clearly reflect the maturity difference. Detection studies have a median sample size of 66 participants compared to 40 for forecasting studies. The largest detection study enrolled 384 patients across eight centers, while the largest forecasting study enrolled 70 patients.

\subsubsection{Commercial Device Status}

Six of 13 studies use commercially available devices. The Empatica E4 research wristband appears in three studies. The Embrace watch appears in one study. The NightWatch armband appears in one study. Fitbit smartwatches appear in one study.

The Embrace watch has FDA 510(k) clearance for seizure detection. NightWatch has CE approval in Europe. However, no device has regulatory approval specifically for seizure forecasting. Forecasting devices would need to establish clinical utility and safety through dedicated trials.

Seven studies use research prototypes without clear commercialization pathways. These prototypes demonstrate technical feasibility but may require substantial engineering for commercial deployment.

\subsubsection{Real-World Deployment}

Six studies include home or ambulatory validation. \textcite{dongTwoLayerEnsembleMethod2022} achieved the most extensive home validation with 788 overnight recordings over three months. \textcite{stirlingForecastingSeizureLikelihood2021} achieved the longest monitoring duration with a mean of 14.6 months per patient.

\textcite{nasseriAmbulatorySeizureForecasting2021} achieved the most rigorous ambulatory validation with concurrent invasive EEG confirmation via the RNS system. However, only six patients were studied.

Hospital settings often fail to capture the diversity of real-world activities that cause false alarms. Exercise, cooking, household chores, and other activities present challenges not encountered in controlled hospital environments.

\subsubsection{Deployment Architecture}

Three studies report on-device or real-time deployment capability, Cloud-based processing also appears in three studies.

\textcite{spahrDeepLearningbasedDetection2025} reported 112 ms inference time suitable for real-time processing without any mention of the specific device achieving it. \textcite{borujenyDetectionEpilepticSeizure2013} reported 316 mW power consumption for server-based processing. \textcite{dongTwoLayerEnsembleMethod2022} deployed on the NightWatch armband with real-time processing.

\textcite{nasseriAmbulatorySeizureForecasting2021} and \textcite{odeDevelopmentEpilepticSeizure2023} both used cloud-based approaches. 

Cloud processing may introduce latency and connectivity issues. Offline processing appears in two studies without clear deployment pathways.

\subsubsection{Barriers to Clinical Adoption}

Several barriers impede clinical translation of wearable seizure monitoring. First, high FPR limits clinical utility. Six of eight detection studies report FPR above acceptable thresholds. \textcite{dongTwoLayerEnsembleMethod2022} reported 0.165/h FPR in home validation, \textcite{wangEpilepticSeizureDetection2025} reported 0.354/h, and \textcite{reintjesECGBasedDetectionEpileptic2025} reported FPR from 0.11/h to 65.62/h depending on method and configuration. High FPR causes alarm fatigue and may lead patients to abandon devices.

Second, limited seizure type coverage restricts applicability. Detection approaches focus on convulsive seizures with motor manifestations captured by accelerometer and gyroscope sensors. Focal aware seizures without motor signs remain undetectable with current wearable approaches. Non-convulsive seizures represent a substantial unmet need that may require invasive EEG monitoring.

Third, small sample sizes limit generalizability. Four studies have sample sizes below 20 patients. \textcite{singhrathoreDevelopmentMultimodalMachine2024} is a single-patient case study. \textcite{borujenyDetectionEpilepticSeizure2013} studied only 3 patients. Small samples increase uncertainty about performance estimates and may not capture the full diversity of seizure presentations.

Fourth, validation setting affects real-world applicability. Seven studies were conducted entirely in hospital or EMU settings. Hospital environments may not capture the diversity of real-world activities that cause false alarms.

Fifth, device burden and battery life limit adherence. Multi-modal systems sampling multiple signals at high frequencies consume substantial power. Long-term monitoring requires daily charging which may reduce adherence over time.
